\label{chap:83-13}

\section{Construction préconstitutive de la réponse bidirectionnelle d'action et de l'opérateur de vide}
\label{p12-83-c13-s1:ch:vacuum}

\subsection{Dérivation de l'échelle d'action à partir de l'état dodécaphasique}

La décennie d'action est déterminée intrinsèquement, sans l'adopter comme
unité héritée. La fonctionnelle déterminantale locale s'applique à la sortie HMT
préalablement construite \(\alpha_{\HMT}\) et définit
\begin{equation}
\Sigma_H=\varphi_{\HMT}\alpha_{\HMT}^{16}.
\label{p12-83-c13-s1:eq:action-determinantal-reader}
\end{equation}
L'exposant \(16\) est l'ordre du conoyau local de la clôture signée ; sa valuation décimale fixe
\begin{equation}
\operatorname{ord}_{10}(\Sigma_H)=-34.
\label{p12-83-c13-s1:eq:action-decade}
\end{equation}

L'origine de l'exposant peut être vérifiée dans ce volume. Sur
l'orbite locale à quatre feuillets agit le bloc de Hadamard
\begin{equation}
H_4=
\begin{pmatrix}
1&1&1&1\\
1&1&-1&-1\\
1&-1&1&-1\\
1&-1&-1&1
\end{pmatrix}.
\label{p12-83-c13-s1:eq:action-hadamard}
\end{equation}

\begin{lemma}[Indice local de la droite d'action]
La forme normale de Smith de \(H_4\) est
\begin{equation}
\operatorname{SNF}(H_4)=\operatorname{diag}(1,2,2,4).
\label{p12-83-c13-s1:eq:action-smith}
\end{equation}
En conséquence,
\[
\operatorname{coker}(H_4)\simeq
\Z/2\Z\oplus\Z/2\Z\oplus\Z/4\Z,
\qquad
|\operatorname{coker}(H_4)|=16.
\]
\end{lemma}

\begin{proof}
Le plus grand commun diviseur des entrées est un ; celui des mineurs
\(2\times2\) non nuls est deux ; celui des mineurs \(3\times3\) est quatre ; et
\(|\det H_4|=16\). Les quotients successifs de ces diviseurs déterminantaux
sont \(1,2,2,4\), qui prouvent \cref{p12-83-c13-s1:eq:action-smith}. Le produit des
invariants donne l'ordre du conoyau.
\end{proof}

La norme de la section \(\alpha_{\HMT}\) sur ces seize feuillets
produit \(\alpha_{\HMT}^{16}\) ; l'auto-échelle \(\varphi_{\HMT}\) agit
une seule fois sur la base, et non une fois par feuillet. D'où
\cref{p12-83-c13-s1:eq:action-determinantal-reader}. Les comparaisons certifiées
\begin{equation}
\alpha_{\HMT}^{16}<10^{-34}
<\varphi_{\HMT}\alpha_{\HMT}^{16}<10^{-33}
\label{p12-83-c13-s1:eq:action-decade-bounds}
\end{equation}
démontrent, sans introduire l'unité SI, que l'ordre décimal est
exactement \(-34\). L'évaluation terminale commence par
\begin{equation}
\Sigma_H
=1.046243838104756580184229208303089\ldots\times10^{-34}.
\label{p12-83-c13-s1:eq:action-sigma-value}
\end{equation}
Le remplacement du conoyau par trois copies avec exposant \(16^3\), ou
l'application réitérée de \(\varphi\) dans chaque feuillet, modifierait la
généalogie. Les deux opérations constituent des critères de réfutation et non
des conventions équivalentes.
La chaîne causale qui est conservée est donc
\[
\APP\to\TRIT\to\TPK\to U_{12}^{\rm sign}
\to\alpha_{\HMT}\to\Sigma_H
\to10^{-34}\U_S\to
\{\hbar_{\rm pre},\hbar_{\rm ret}\}.
\]
Ce volume ne redérive pas \(\alpha_{\HMT}\), mais développe la géométrie
métrologique postérieure à la sélection du type et de la décennie par
\(\Sigma_H\). L'entier \(-34\) n'est ni une mantisse ni un ajustement au Système
international, mais la valuation terminale de la fonctionnelle locale.

\subsection{Sections orientées de la droite d'action}

Les réalisations antérieure et postérieure au retour constituent deux sections
d'une même droite d'action :

\begin{equation}
\hbar_{\rm pre}=H_5\,10^{-34}\U_S,\qquad
\hbar_{\rm ret}=\eta_{\rm ret}\,10^{-34}\U_S.
\label{p12-83-c13-s1:eq:action-twoway}
\end{equation}

La décennie commune détermine la droite, tandis que l'information
orientée s'exprime par

\begin{equation}
R_{\rm act}=\frac{\hbar_{\rm pre}}{\hbar_{\rm ret}}
=\frac{H_5}{\eta_{\rm ret}},
\qquad
\xi_{\rm act}=\frac12\log R_{\rm act}.
\label{p12-83-c13-s1:eq:action-rapidity}
\end{equation}

La première quantité fournit une coordonnée multiplicative et la seconde
une coordonnée additive de type rapidité. Elles ne représentent pas deux constantes de
Planck indépendantes, mais les deux orientations d'une même section
par rapport au retour.

Si \(h_a=2\pi\hbar_a\), \(a\in\{\mathrm{pre},\mathrm{ret}\}\), une grandeur proportionnelle à \(h_a^{1/2}\) varie comme \(R_{\rm act}^{1/2}\), une grandeur proportionnelle à \(h_a^{-1/2}\) varie comme \(R_{\rm act}^{-1/2}\), et un quotient où l'action s'annule est un invariant bidirectionnel.

\subsection{Décomposition spectrale bicouche de l'opérateur angulaire}

Soit \(R=R^*=R^{-1}\) l'involution active et

\[
P_\pm=\frac{I\pm R}{2},
\qquad
T=q_+P_+ +q_-P_-,
\qquad 0<q_+<q_-<1.
\]

Avec
\[
x=A_{\rm rad}=\frac{\pi A}{180},\qquad
y=C^*_{\rm rad}=\frac{\pi C^*}{180},
\]
les deux canaux sont construits avant d'évaluer le vide :
\begin{equation}
q_+=e^{-x-y},\qquad q_-=e^{-x+y}.
\label{p12-83-c13-s1:eq:vacuum-angular-channels}
\end{equation}
L'échange de feuillets inverse \(C^*\), c'est-à-dire \(y\mapsto-y\), et permute \(q_+\leftrightarrow q_-\). Le couple est ainsi la lecture spectrale d'un même antécédent angulaire, et non deux paramètres ajustés séparément.

\begin{definition}[Réponse constitutive]
\begin{equation}
\mathcal V_{90,120}=(I-T^{90})(I-T^{120})^{-1}.
\label{p12-83-c13-s1:eq:vacuum-operator}
\end{equation}
\end{definition}

Comme \(\|T\|<1\), le dénominateur est positif et inversible. Le calcul fonctionnel donne

\begin{equation}
\mathcal V_{90,120}=r_+P_+ +r_-P_-,
\qquad
r_\pm=\frac{1-q_\pm^{90}}{1-q_\pm^{120}}.
\label{p12-83-c13-s1:eq:vacuum-eigenvalues}
\end{equation}

L'affirmation d'inversibilité mérite d'être explicitée. Pour chaque feuillet,
\(0<q_\pm^{120}<1\) ; par conséquent
\[
I-T^{120}=(1-q_+^{120})P_++(1-q_-^{120})P_-
\]
est strictement positif et
\[
(I-T^{120})^{-1}
=\frac{P_+}{1-q_+^{120}}+
 \frac{P_-}{1-q_-^{120}}.
\]
La réponse \(\mathcal V_{90,120}\) ne s'obtient pas en résolvant quatre
équations scalaires séparées : elle est le résultat du calcul fonctionnel d'une
seule contraction positive. Cette observation empêche que perméabilité,
permittivité, impédance et propagation perdent l'antécédent qu'elles partagent.

\begin{theorem}[Détermination spectrale conjointe des quatre relations constitutives]
Les quatre caractères constitutifs
\begin{equation}
\widehat\varepsilon=r_+^2,\qquad
\widehat\mu=r_-^2,\qquad
\widehat Z=\frac{r_-}{r_+},\qquad
\widehat c=\frac1{r_+r_-}
\label{p12-83-c13-s1:eq:vacuum-four}
\end{equation}
sont des fonctions spectrales de \(\mathcal V_{90,120}\). Ils satisfont
\begin{equation}
\widehat\mu\widehat\varepsilon=\widehat c^{-2},
\qquad
\frac{\widehat\mu}{\widehat\varepsilon}=\widehat Z^2.
\label{p12-83-c13-s1:eq:vacuum-relations}
\end{equation}
\end{theorem}

\begin{proof}
Les identités s'obtiennent en substituant \cref{p12-83-c13-s1:eq:vacuum-four}. Aucune carte dimensionnelle ni aucune valeur extérieure n'intervient.
\end{proof}

\begin{corollary}[Reconstruction des deux feuillets]
Les quatre caractères ne contiennent pas quatre degrés de liberté. N'importe lequel
des couples \((\widehat\mu,\widehat\varepsilon)\),
\((\widehat Z,\widehat c)\) ou \((\mathcal E,\mathcal B)\) reconstruit les
deux valeurs propres positives :
\begin{align}
r_-&=\sqrt{\widehat\mu}
=\sqrt{\widehat Z/\widehat c},
&r_+&=\sqrt{\widehat\varepsilon}
=\frac1{\sqrt{\widehat Z\widehat c}},
\label{p12-83-c13-s1:eq:vacuum-reconstruct-r}\\
\log r_-&=\frac{\mathcal E+\mathcal B}{2},
&\log r_+&=\frac{\mathcal E-\mathcal B}{2}.
\label{p12-83-c13-s1:eq:vacuum-reconstruct-log}
\end{align}
En particulier, l'orientation se perd si l'on conserve uniquement le produit,
mais se retrouve en ajoutant le quotient.
\end{corollary}

\begin{proof}
La première ligne est la racine positive des identités
\cref{p12-83-c13-s1:eq:vacuum-four,p12-83-c13-s1:eq:vacuum-relations}. La seconde inverse le système
linéaire
\(\mathcal E=\log r_-+\log r_+\),
\(\mathcal B=\log r_--\log r_+\).
\end{proof}

Ce corollaire exprime une propriété centrale de HMT : produit et quotient sont
des fonctionnelles complémentaires d'un même état. Le produit détermine la
propagation commune, tandis que le quotient conserve l'orientation
relative des feuillets intérieur et extérieur.

Dans le certificat V2,

\begin{align*}
r_+&=0.9999996285482493631786952281\ldots,\\
r_-&=0.9997241371895183641315165853\ldots,\\
\widehat\varepsilon&=0.9999992570966367027604416155\ldots,\\
\widehat\mu&=0.9994483504793269350899815559\ldots,\\
\widehat Z&=0.9997245085389372154555543466\ldots,\\
\widehat c&=1.0002763104861575261063862689\ldots.
\end{align*}

Ces numéraux sont la reconnaissance terminale de \cref{p12-83-c13-s1:eq:vacuum-operator} ; ils ne définissent pas ses exposants \(90,120\).

\subsection{Factorisation harmonique des secteurs \texorpdfstring{\(90\)}{90} et \texorpdfstring{\(120\)}{120}}

Le semi-groupe harmonique est

\[
\langle90,120\rangle=30\langle3,4\rangle.
\]

Si \(s=q^{30}\), alors

\begin{equation}
\frac{1-q^{90}}{1-q^{120}}
=\frac{1-s^3}{1-s^4}
=\frac{1+s+s^2}{1+s+s^2+s^3}.
\label{p12-83-c13-s1:eq:three-four-completion}
\end{equation}

L'identité sépare également la réponse de son défaut. Si
\[
F_{3\to4}(s)=\frac{1+s+s^2}{1+s+s^2+s^3},
\qquad 0<s<1,
\]
alors
\begin{equation}
d_{3\to4}(s):=1-F_{3\to4}(s)
=\frac{s^3}{1+s+s^2+s^3}>0.
\label{p12-83-c13-s1:eq:vacuum-defect34}
\end{equation}
Le terme supplémentaire est conservé exactement comme défaut
positif. En outre,
\begin{equation}
F_{3\to4}'(s)
=-\frac{s^2(3+2s+s^2)}{(1+s+s^2+s^3)^2}<0.
\label{p12-83-c13-s1:eq:vacuum-monotonic34}
\end{equation}
Par conséquent, l'ordre des canaux angulaires détermine univoquement
l'ordre des réponses constitutives. L'attribution de
\(\widehat\mu\) et \(\widehat\varepsilon\) aux feuillets respectifs s'obtient
analytiquement, sans recourir à leurs développements décimaux.

Chaque feuillet du vide constitue la complétion exacte de trois à quatre
termes d'un même mode élémentaire \(30\). En conséquence, perméabilité
et permittivité s'interprètent comme deux fonctionnelles quadratiques d'une même
complétion évaluée sur des feuillets conjugués, et non comme des paramètres
constitutifs indépendants.

Comme la fonction \((1-q^{90})/(1-q^{120})\) décroît dans \(0<q<1\) et \(q_+<q_-\), on obtient

\[
0<r_-<r_+<1,\qquad
0<\widehat\mu<\widehat\varepsilon<1,\qquad
0<\widehat Z<1<\widehat c.
\]

Les limites de la fonctionnelle déterminent sa géométrie. Lorsque \(q\to0^+\),
\(F(q)\to1\) : les quatre lectures s'approchent de l'unité et le défaut de
mémoire disparaît. Lorsque \(q\to1^-\),
\[
F(q)\longrightarrow\frac{90}{120}=\frac34.
\]
La réponse prend donc ses valeurs dans l'intervalle positif \((3/4,1)\). La
borne \(3/4\) n'est pas une constante ajustée : elle est le quotient des deux
longueurs de canal et la réalisation continue de la complétion
trois--à--quatre.

\subsection{Involution des feuillets et conservation de l'orientation}

L'échange de feuillets \(C^*\mapsto-C^*\) agit par

\begin{equation}
\widehat\mu\longleftrightarrow\widehat\varepsilon,\qquad
\widehat Z\longmapsto\widehat Z^{-1},
\qquad
\widehat c\longmapsto\widehat c.
\label{p12-83-c13-s1:eq:vacuum-sheet}
\end{equation}

La propagation commune est paire ; l'impédance conserve l'orientation.
Cette distinction réapparaîtra dans Barbero et CKM : toutes les fonctionnelles
associées au changement de feuillet n'ont pas la même parité.

\subsection{Caractères logarithmiques et compatibilité traciale}

Définissons

\[
\mathcal E=\log(r_+r_-),\qquad
\mathcal B=\log(r_-/r_+).
\]

Avec la trace normalisée \(\tau=\Tr/2\),

\begin{equation}
2\tau(\log\mathcal V)=\mathcal E,
\qquad
-2\tau(R\log\mathcal V)=\mathcal B.
\label{p12-83-c13-s1:eq:vacuum-traces}
\end{equation}

Le couple \((\mathcal E,\mathcal B)\) constitue un système de coordonnées linéaire du logarithme
opératoriel :
\begin{equation}
\log\mathcal V
=\frac{\mathcal E}{2}I-\frac{\mathcal B}{2}R.
\label{p12-83-c13-s1:eq:vacuum-log-decomposition}
\end{equation}
La composante scalaire est invariante sous l'échange de feuillets ; la
composante orientée change de signe. Cette décomposition démontre
directement la parité de la propagation et le caractère orienté de
l'impédance.

L'amplification nonadique

\[
\mathcal V_m=\mathcal V\otimes I_{9^m}
\]

est compatible avec les espérances \(E_m=\mathrm{id}\otimes\tau_9\) :

\[
E_m(\mathcal V_{m+1})=\mathcal V_m.
\]

Par conséquent, \(\mathcal E\) et \(\mathcal B\) ne sont pas des particularités d'une matrice \(2\times2\) ; ils subsistent dans la tour UHF et dans sa clôture traciale \(II_1\).

Plus précisément, pour tout polynôme \(p\) puis, par continuité,
pour toute fonction continue sur le spectre,
\begin{equation}
E_m\bigl(p(\mathcal V_{m+1})\bigr)=p(\mathcal V_m).
\label{p12-83-c13-s1:eq:vacuum-functional-refinement}
\end{equation}
La compatibilité ne préserve pas seulement deux nombres : elle préserve toute
l'algèbre fonctionnelle engendrée par la réponse. En particulier, elle conserve ses
projecteurs spectraux, ses puissances, son logarithme et les formes quadratiques
qui mesurent les défauts \cref{p12-83-c13-s1:eq:vacuum-defect34}.

\subsection{Réalisation dimensionnelle des relations constitutives}

La mécanique dimensionnelle détermine

\begin{align}
Z_{\rm vac}&=\widehat Z\,u_Su_Q^{-2},\\
c_{\rm vac}&=\widehat c\,u_C,\\
\mu_{\rm vac}&=\widehat\mu\,u_Su_C^{-1}u_Q^{-2},\\
\varepsilon_{\rm vac}&=\widehat\varepsilon\,u_S^{-1}u_C^{-1}u_Q^2.
\label{p12-83-c13-s1:eq:vacuum-sections}
\end{align}

Les identités constitutives se maintiennent dans les droites :

\begin{equation}
\mu_{\rm vac}\varepsilon_{\rm vac}=c_{\rm vac}^{-2},
\qquad
\mu_{\rm vac}/\varepsilon_{\rm vac}=Z_{\rm vac}^2.
\end{equation}

La carte SI est postérieure. Il n'est pas nécessaire de déclarer \(c\) primitif et de résoudre ensuite \(\mu\varepsilon=c^{-2}\) : les trois objets descendent conjointement de \(\mathcal V\) et des droites déclarées.

La séparation des niveaux s'exprime par le diagramme
\[
\begin{array}{c}
(A,C^*)\longrightarrow T\longrightarrow\mathcal V
\longrightarrow
(\widehat\mu,\widehat\varepsilon,\widehat Z,\widehat c)
\\[2mm]
\Big\downarrow\text{ cartes de droite}\hspace{36mm}
\Big\downarrow\text{ évaluation finale}
\\[2mm]
(u_S,u_Q,u_C)\longrightarrow
(\mu_{\rm vac},\varepsilon_{\rm vac},Z_{\rm vac},c_{\rm vac})
\longrightarrow\text{ coordonnées SI}.
\end{array}
\]
La première ligne est adimensionnelle et opératorielle ; la deuxième détermine les types de
grandeur ; la troisième sélectionne une carte métrologique. Modifier la carte finale
ne modifie ni les valeurs propres ni les identités constitutives. En ce
sens précis, l'unification mathématique--physique consiste à distinguer
l'invariant généalogique de son système de coordonnées, et non à supprimer la
structure dimensionnelle.

\begin{proposition}[Minimalité spectrale]
Soit \(W\) un opérateur positif à deux feuillets de spectre simple
\(\{r_+,r_-\}\). Si quatre caractères positifs satisfont
\(\mu\varepsilon=c^{-2}\) et \(\mu/\varepsilon=Z^2\), si la propagation
commune est déterminée par \(c^{-1}=\det W=r_-r_+\), et si son orientation
fixe \(Z=r_-/r_+\), alors nécessairement
\[
(\mu,\varepsilon,Z,c)
=\left(r_-^2,r_+^2,\frac{r_-}{r_+},\frac1{r_-r_+}\right).
\]
Par conséquent, \cref{p12-83-c13-s1:eq:vacuum-four} n'est pas une paramétrisation parmi d'autres :
c'est la réalisation positive minimale compatible avec le produit, le quotient et
l'orientation.
\end{proposition}

\begin{proof}
De \(Z=r_-/r_+\) et \(c^{-1}=\det W=r_-r_+\), on obtient de manière unique les
racines positives. Les deux identités constitutives imposent ensuite
\(\mu=r_-^2\) et \(\varepsilon=r_+^2\).
\end{proof}

\subsection{Déterminant angulaire et relation électrofaible de Weinberg}

Le déterminant de \(T\) est

\begin{equation}
\det T=q_+q_-=D_A.
\label{p12-83-c13-s1:eq:detT}
\end{equation}

Le secteur constitutif du vide évalue \(F(T)^2\), où
\(F(z)=(1-z^{90})/(1-z^{120})\), tandis que le secteur angulaire
électrofaible évalue \(\det T\). Il s'agit de deux fonctions d'un
antécédent commun, et non d'une factorisation horizontale entre constantes.
Cette distinction est développée dans le \cref{ch:ew}.

\begin{falsifier}
Le bloc est réfuté si l'une quelconque de ces identités exactes est violée :
positivité de \(I-T^{120}\), égalité
\cref{p12-83-c13-s1:eq:three-four-completion}, relations \cref{p12-83-c13-s1:eq:vacuum-relations} ou
compatibilité \(E_m(\mathcal V_{m+1})=\mathcal V_m\). Une coïncidence SI
ne peut pas compenser une défaillance opératorielle.
\end{falsifier}

\begin{statusbox}
Opérateur, spectre, complétion \(3\to4\), parité, amplification et
sections constitutives forment un théorème interne exact. Leur réunion en un
seul opérateur est établie dans cette formulation ; les données angulaires et
les droites, dérivées précédemment dans HMT, constituent la structure de départ
explicite.
\end{statusbox}


\section{Constantes quantiques--électriques, échelle thermique et système de Planck}
\label{p12-83-c13-s2:ch:metrology}

\subsection{Dérivation de la section de charge à partir du vide}

\(Z_{\rm vac}\), \(h_a=2\pi\hbar_a\) et \(\alpha_{\HMT}\) étant fixés, la charge du feuillet \(a\in\{\mathrm{pre},\mathrm{ret}\}\) est

\begin{equation}
e_a=\sqrt{\frac{2\alpha_{\HMT}h_a}{Z_{\rm vac}}}.
\label{p12-83-c13-s2:eq:charge}
\end{equation}

La section \(e_a\) n'est pas introduite au moyen d'une coordonnée du Système
international destinée à reconstruire \(\alpha\). L'ordre causal procède de
\(\alpha_{\HMT}\), de l'action et du vide vers la section de charge.

\subsection{Constantes de von Klitzing et de Josephson, conductance et flux magnétique}

De \cref{p12-83-c13-s2:eq:charge} se déduisent

\begin{align}
R_K&=\frac{h_a}{e_a^2}=\frac{Z_{\rm vac}}{2\alpha_{\HMT}},
\label{p12-83-c13-s2:eq:rk}\\
G_0&=\frac{2e_a^2}{h_a}=\frac{4\alpha_{\HMT}}{Z_{\rm vac}},
\label{p12-83-c13-s2:eq:g0}\\
\Phi_{0,a}&=\frac{h_a}{2e_a}
=\sqrt{\frac{h_aZ_{\rm vac}}{8\alpha_{\HMT}}},
\label{p12-83-c13-s2:eq:phi0}\\
K_{J,a}&=\frac{2e_a}{h_a}=\Phi_{0,a}^{-1}.
\label{p12-83-c13-s2:eq:kj}
\end{align}

En conséquence,

\begin{equation}
R_KG_0=2,\qquad G_0Z_{\rm vac}=4\alpha_{\HMT},
\qquad K_{J,a}\Phi_{0,a}=1.
\label{p12-83-c13-s2:eq:quantum-electrical-identities}
\end{equation}

La covariance bidirectionnelle prend la forme

\[
R_K^{\rm pre}=R_K^{\rm ret},\qquad
G_0^{\rm pre}=G_0^{\rm ret},
\]

\[
\frac{\Phi_{0,\rm pre}}{\Phi_{0,\rm ret}}=R_{\rm act}^{1/2},
\qquad
\frac{K_{J,\rm pre}}{K_{J,\rm ret}}=R_{\rm act}^{-1/2}.
\]

Les relations de résistance et de conductance sont indépendantes de
l'orientation de l'action. Le couple flux--fréquence conserve cette
orientation et transforme réciproquement ses deux composantes.

\subsection{Relation thermique intrinsèque de l'échelle \texorpdfstring{\(108\)}{108}}

La structure chronologique \(108\) fixe

\begin{equation}
k_B\Theta_{\rm clk}\log3
=\frac{h_{\rm int}}{108t_0}
=\frac{2\pi\hbar_{\rm int}}{108t_0}.
\label{p12-83-c13-s2:eq:boltzmann-clock}
\end{equation}

Cette égalité ne détermine pas séparément une coordonnée numérique de \(k_B\) et une autre de \(\Theta_{\rm clk}\) sans normalisation supplémentaire. Elle détermine bien leur produit comme section d'énergie. Le facteur \(\log3\) possède trois générateurs compatibles :

\begin{equation}
\log3=C_1+C_2-2C_0
=\frac{S_{\TRIT}}{k_B}
=\frac{E_P}{k_B\Theta_{\rm clk}}.
\label{p12-83-c13-s2:eq:log3-triple}
\end{equation}

Les trois égalités correspondent aux réalisations résiduelle, informationnelle et
thermodynamique d'un même scalaire. La coïncidence de valeur n'identifie pas
leurs domaines respectifs.

\subsection{Système de Planck et covariance bidirectionnelle}

Le sélecteur CT108 construit

\begin{equation}
\theta_{108}=\left(\frac{54}{\pi}\right)^2,
\qquad
\ell_P=\frac{54}{\pi}\ell_0,
\qquad
t_P=\frac{54}{\pi}t_0,
\label{p12-83-c13-s2:eq:planck-length-time}
\end{equation}

et

\begin{equation}
E_P=\frac{\hbar_{\rm int}}{t_P}
=\frac{\pi\hbar_{\rm int}}{54t_0}.
\label{p12-83-c13-s2:eq:planck-energy}
\end{equation}

En comparant \cref{p12-83-c13-s2:eq:boltzmann-clock,p12-83-c13-s2:eq:planck-energy}, on obtient

\begin{equation}
\boxed{E_P=k_B\Theta_{\rm clk}\log3.}
\label{p12-83-c13-s2:eq:planck-trit}
\end{equation}

La relation ne dépend pas d'une coïncidence entre unités : les deux membres
proviennent de la même structure chronologique \(108\).

Lorsque la constante gravitationnelle du \cref{ch:gravity} est introduite, la famille de Planck satisfait

\begin{align}
\ell_P^2&=\frac{G\hbar}{c^3}=\theta_{108}\ell_0^2,
\label{p12-83-c13-s2:eq:planck-area}\\
F_P&=\frac{c^4}{G},
\qquad
P_P=\frac{c^5}{G},
\label{p12-83-c13-s2:eq:planck-force-power}\\
\rho_P&=\frac{\hbar}{\theta_{108}^2c^5t_0^4}.
\label{p12-83-c13-s2:eq:planck-density}
\end{align}

Bien que \(G\) et \(\hbar\) se transforment réciproquement entre les sections
antérieure et postérieure au retour, leur produit dans \(\ell_P^2\) conserve la
géométrie. Cette invariance détermine l'aire comme interface entre les
réalisations géométrique, informationnelle et d'intrication.

\subsection{Formulations spectrales de la loi de déplacement de Wien}

Pour la densité spectrale par longueur d'onde, le maximum satisfait

\[
5(1-\ee^{-x_5})=x_5,
\qquad
x_5=5+W_0(-5\ee^{-5}).
\]

Avec \cref{p12-83-c13-s2:eq:boltzmann-clock},

\begin{equation}
\lambda_{\max}(\Theta_{\rm clk})
=\frac{108\log3}{x_5}\,\ell_0.
\label{p12-83-c13-s2:eq:wien-lambda}
\end{equation}

Pour la densité par fréquence,

\[
3(1-\ee^{-x_3})=x_3,
\qquad
x_3=3+W_0(-3\ee^{-3}),
\]

et

\begin{equation}
\nu_{\max}(\Theta_{\rm clk})
=\frac{x_3}{108t_0\log3}.
\label{p12-83-c13-s2:eq:wien-frequency}
\end{equation}

Les maxima ne sont pas réciproques, puisque les densités spectrales se
transforment avec le jacobien correspondant. Le quotient \(x_5/x_3\)
appartient à la structure de la paramétrisation et ne représente pas une
incohérence.

\subsection{Loi de Stefan--Boltzmann et thermodynamique du gaz de photons}

Évaluer la loi de rayonnement en \(\Theta_{\rm clk}\) produit le flux intrinsèque

\begin{equation}
j_\star(\Theta_{\rm clk})
=\frac{2\pi^5}{15\,108^4(\log3)^4}\,
\frac{h_{\rm int}}{\ell_0^2t_0^2},
\label{p12-83-c13-s2:eq:stefan-hmt}
\end{equation}

dans la normalisation déclarée. La constante de Stefan--Boltzmann n'est pas insérée comme primitive : elle résulte de l'intégration du spectre bosonique une fois l'action, la propagation et la température fixées.
Les lois spectrales employées sont celles de Planck et de Wien \cite{rm:Planck,rm:Wien} ; HMT fixe ici la section sur laquelle elles sont évaluées, et non leurs constantes par comparaison rétrospective.

En unités de l'échelle de Planck,

\begin{align}
n_\gamma\ell_P^3
&=\frac{2\zeta(3)}{\pi^2\log^3 3},
\label{p12-83-c13-s2:eq:photon-number}\\
\frac{u_\gamma\ell_P^3}{E_P}
&=\frac{\pi^2}{15\log^4 3},
\label{p12-83-c13-s2:eq:photon-energy}\\
\frac{s_\gamma\ell_P^3}{k_B}
&=\frac{4\pi^2}{45\log^3 3}.
\label{p12-83-c13-s2:eq:photon-entropy}
\end{align}

\(\zeta(3)\) représente le moment bosonique d'ordre trois, tandis que
\(\log3\) détermine l'énergie de décision TRIT. L'équation conserve la
distinction entre les deux invariants.

La conductance thermique quantique s'exprime comme

\begin{equation}
\frac{\kappa_Qt_0}{k_B}=\frac{\pi^2}{324\log3}.
\label{p12-83-c13-s2:eq:thermal-conductance}
\end{equation}

\subsection{Confluence métrologique ultérieure de l'action, de la charge, de la conductance et du flux magnétique}

\begin{longtable}{@{}p{0.18\textwidth}p{0.27\textwidth}p{0.45\textwidth}@{}}
\toprule
Objet & Équation interne & Significado estructural\\
\midrule
\endhead
\(R_K\) & \(Z_{\rm vac}/(2\alpha)\) & résistance quantique invariante bidirectionnelle\\
\(G_0\) & \(4\alpha/Z_{\rm vac}\) & conductancia complementaria\\
\(\Phi_0\) & \(\sqrt{hZ/(8\alpha)}\) & flux qui conserve l'orientation de l'action\\
\(K_J\) & \(\Phi_0^{-1}\) & relation réciproque entre fréquence et flux\\
\(E_P\) & \(k_B\Theta\log3\) & énergie du cycle et contenu informationnel TRIT\\
Wien & \(108\log3/x_5\) & extremum spectral à l'échelle chronologique\\
Apéry photonique & \(2\zeta(3)/(\pi^2\log^3 3)\) & densité de nombre bosonique\\
\bottomrule
\end{longtable}

\begin{falsifier}
Le bloc est réfuté si un changement de variable spectrale conserve
indûment le même extremum de Wien, si \(\zeta(3)\) remplace
\(\log3\), ou si des valeurs séparées sont attribuées à \(k_B\) et \(\Theta\) sans
spécifier de carte. Les identités
\cref{p12-83-c13-s2:eq:quantum-electrical-identities,p12-83-c13-s2:eq:planck-trit} fournissent des contrôles
algébriques directs.
\end{falsifier}

\begin{statusbox}
Quanta électriques et confluence Planck--TRIT : \status{Théorème interne exact}. Wien, Stefan--Boltzmann et gaz photonique : \status{Composition mathématique exacte} avec le continu bosonique déclaré. La carte SI demeure un \status{Contraste métrologique externe}.
\end{statusbox}


\section{Réalisation analytique graduée de la microfibre quintuple}
\label{p12-83-c13-s3:chap:realizacion-analitica-contraangulo}
\index{coordonnée angulaire conjuguée \(C^*\)}
\index{microfibre de pi@microfibre de \(\pi\)!caractère analytique}
\index{caractère déterminantal}

Le catalogue fini et une évaluation analytique répondent à des questions de types
distincts. Le catalogue détermine quelles régions existent, lesquelles partagent un
retour et quelle information chacune conserve. Une évaluation analytique doit
en outre décider comment la longueur d'un mot se transforme en degré, comment
une branche se prolonge après son retour et avec quel caractère scalaire cette
prolongation se lit. Ce chapitre construit cette seconde opération de manière
explicite.

La construction part de quatre données déjà obtenues dans les chapitres
précédents : la microfibre à cinq régions de \(\pi\), la longueur six de
chaque mot régional, la clôture dodécaphasique de \(\alpha\) et le transport
bicouche associé à la coordonnée commune
\(A=10^3\alpha\). Aucune valeur de la coordonnée angulaire conjuguée \(C_*^{\rm HMT}\), de la composante réduite
d'action, de la fonctionnelle de Barbero--Immirzi ou d'une constante métrologique n'est
employée dans l'évaluation. Le résultat est un caractère régional convergent,
une récurrence exacte et une phase orientée dont la carte sexagésimale reproduit
la coordonnée angulaire conjuguée publiée.

La séparation entre la partie finie et la partie analytique est essentielle. L'entier
\(169\) sera un théorème du catalogue. La série qui le prolonge sera le
théorème d'un lecteur analytique dont les règles sont déclarées avant de l'évaluer.
Ainsi, les prémisses supplémentaires ne restent pas dissimulées dans une coïncidence
décimale.

\subsection{Le retour persistant des \(5\) régions}

Soit
\[
\mathcal F_\pi=\Psi^{-1}(010211)\subset\mathcal U_6
\]
la microfibre relative au catalogue. Rappelons ses cinq éléments, maintenant séparés en bloc initial et bloc terminal :
\[
\begin{aligned}
 &(501,614,498\mid169,272,272),\\
 &(810,923,870\mid169,272,272),\\
 &(810,983,810\mid169,272,272),\\
 &(870,923,810\mid169,272,272),\\
 &(870,923,810\mid418,674,521).
\end{aligned}
\]
Nous définissons la projection de retour
\[
p_{\rm ret}:\mathbb Z^6\longrightarrow\mathbb Z^3,
\qquad
p_{\rm ret}(u_1,\ldots,u_6)=(u_4,u_5,u_6).
\]

\begin{proposition}[Bloc terminal persistant]
\label{p12-83-c13-s3:prop:bloque-terminal-persistente}
Le multiensemble \(p_{\rm ret}(\mathcal F_\pi)\) possède un unique élément de
multiplicité maximale :
\[
b_\pi=(169,272,272),
\qquad
\operatorname{mult}_{\mathcal F_\pi}(b_\pi)=4.
\]
L'unique bloc terminal restant est
\[
b_\pi^-=(418,674,521)
\]
et a une multiplicité de un.
\end{proposition}

\begin{proof}
L'affirmation s'obtient par énumération exacte des cinq lignes du catalogue qui satisfont \(\Psi(U)=010211\). Quatre lignes possèdent le bloc terminal \((169,272,272)\) ; la cinquième, qui porte le retour muté, possède \((418,674,521)\). La multiplicité maximale est quatre et n'est atteinte qu'une seule fois.
\end{proof}

La différence orientée entre le retour muté et le retour persistant est
\[
 \omega_\pi=b_\pi^- - b_\pi=(249,402,249),
 \qquad
 \langle(1,1,1),\omega_\pi\rangle=900.
\]
Cette identité sépare deux objets qu'il ne faut pas confondre. L'entier
\(169\) est une coordonnée d'une donnée terminale, c'est-à-dire une donnée de degré
zéro. Le vecteur \(\omega_\pi\) peut être lu comme une cochaîne de degré un
après orientation de l'arête qui relie les deux blocs. L'appeler cocycle exigerait,
en outre, de fixer un complexe de cochaînes et de vérifier que son cobord s'annule.
La prolongation analytique construite plus loin ne présuppose pas cette
identification.

La sélection de \(169\) n'exige pas de privilégier la première position du bloc.
Le groupe symétrique \(\mathfrak S_3\) agit en permutant ses trois coordonnées.
Le stabilisateur de \(b_\pi\) échange les deux occurrences de \(272\) et
fixe l'unique coordonnée de multiplicité un. Notons
\(\operatorname{sing}(b)\) cette coordonnée lorsque le bloc est de type
\((r,s,s)\), avec \(r\ne s\). De manière équivalente,
\[
\operatorname{sing}(b)
=\sum_{j=1}^3b_j-2\operatorname{moda}(b).
\]
En particulier,
\[
\boxed{\operatorname{sing}(b_\pi)=169.}
\]

\begin{definition}[Sélecteur résiduel de la microfibre]
\label{p12-83-c13-s3:def:selector-residual-pi}
Le sélecteur \(\rho_\pi\) applique successivement deux opérations :
\begin{enumerate}
  \item sélectionne le bloc terminal de multiplicité maximale au sein de \(p_{\rm ret}(\mathcal F_\pi)\) ;
  \item retient la coordonnée fixe sous l'action du stabilisateur de ce bloc.
\end{enumerate}
Par la \cref{p12-83-c13-s3:prop:bloque-terminal-persistente},
\[
\boxed{\rho_\pi=169.}
\]
\end{definition}

\begin{theorem}[Unicité équivariante du retour]
\label{p12-83-c13-s3:thm:unicidad-retorno-169}
Parmi les sélecteurs qui
\begin{enumerate}
  \item sont invariants sous les permutations des cinq régions ;
  \item sont équivariants sous les permutations des trois coordonnées
  terminales ;
  \item sélectionnent le bloc de multiplicité maximale ;
  \item retiennent une coordonnée fixe sous l'action de son stabilisateur,
\end{enumerate}
la valeur de retour est déterminée de manière unique et vaut \(169\).
\end{theorem}

\begin{proof}
L'invariance sous \(\mathfrak S_5\) empêche d'utiliser l'ordre des lignes. La
troisième condition sélectionne \(b_\pi\), qui est l'unique mode d'après la
\cref{p12-83-c13-s3:prop:bloque-terminal-persistente}. Son stabilisateur dans
\(\mathfrak S_3\) possède deux orbites sur les coordonnées : la paire de
valeurs \(272\) et la valeur singleton \(169\). La quatrième condition sélectionne la
seconde orbite. L'équivariance conserve sa valeur lorsque sa position est déplacée.
\end{proof}

Cet argument améliore la lecture positionnelle du retour. Le \(169\) n'est pas
« la première coordonnée » d'une ligne choisie : c'est un invariant du
multiensemble régional et de l'action de ses symétries de réordonnancement.
Il n'est pas non plus sélectionné pour sa rareté globale. Dans les \(468\) émissions du
catalogue, \(169\) apparaît \(84\) fois, uniformément réparti avec
\(14\) occurrences dans chacune des six positions. L'unicité du
\cref{p12-83-c13-s3:thm:unicidad-retorno-169} est donc relationnelle et régionale : elle réside
dans la combinaison de la microfibre, de la persistance modale et du stabilisateur, non dans la
fréquence isolée de l'entier.

\subsection{Longueur cylindrique et degré analytique}

Soit \(\mathcal P\) le module libre engendré par les mots régionaux et
soit \(F^n\mathcal P\) sa filtration par longueur. La graduation associée
enregistre le premier niveau auquel apparaît chaque mot. Pour un paramètre
\(z\), le caractère ordinaire de longueur est
\[
\chi_z([w])=z^{|w|}.
\]
La multiplicativité
\[
\chi_z([uv])=\chi_z([u])\chi_z([v])
\]
exprime que la concaténation additionne les longueurs. Lors de l'évaluation en
\(z=\alpha\), un mot de longueur \(n\) reçoit le degré \(\alpha^n\).

\begin{definition}[Compatibilité de degré]
\label{p12-83-c13-s3:def:compatibilidad-grado}
Le lecteur analytique régional est compatible avec la filtration cylindrique
lorsqu'il envoie la composante homogène de longueur \(n\) sur le degré analytique
\(n\) :
\[
\operatorname{gr}_n\mathcal P\longrightarrow
\alpha^n\mathbb R.
\]
\end{definition}

Les cinq régions appartiennent à \(\mathcal U_6\) ; c'est pourquoi la donnée de
retour apparaît en degré six :
\[
\rho_\pi[U_6]\longmapsto169\alpha^6.
\]
Ce six est la longueur du mot régional. Ce n'est pas la longueur des
marches fermées du \cref{chap:cinco-ciclos-pi}. Ces marches parcourent à
la fois les quatre ancrages
\(U_\pi,U_e,U_\varphi,U_{\rm APP}\) et possèdent une longueur minimale de huit. Un
mot, une arête du graphe de blocs et une marche fermée sont des objets de
types distincts ; la graduation utilisée ici appartient au premier.

\subsection{Le caractère déterminantal du transport commun}

La fermeture dodécaphasique fournit \(\alpha\). La complétion des trois
couples nonadiques fournit l'échelle \(10^3\), de sorte que
\[
A=10^3\alpha.
\]
La carte angulaire archimédienne étant choisie, la coordonnée commune en radians est
\[
x=\frac{\pi A}{180}=\frac{50\pi}{9}\alpha.
\]
Soit \(R\) une involution réelle de trace nulle en dimension deux,
\[
R^2=I_2,
\qquad
\operatorname{tr}R=0,
\]
et considérons le transport formel bicouche
\[
\mathcal T(x,y)=\exp(-xI_2-yR).
\]
La variable orientée \(y\) n'a pas encore été évaluée. Cependant, l'action
de \(\mathcal T\) sur la droite déterminante en est indépendante :
\[
\begin{aligned}
\det\mathcal T(x,y)
&=\exp\!\left(\operatorname{tr}(-xI_2-yR)\right)\\
&=e^{-2x}.
\end{aligned}
\]
On définit
\[
\boxed{
D_A=e^{-2x}=e^{-100\pi\alpha/9}.}
\]
L'indice rappelle qu'il s'agit du déterminant de la contraction commune associée à \(A\), et non d'une fonction de \(C_*^{\rm HMT}\).

\begin{proposition}[Caractère de la droite déterminante]
\label{p12-83-c13-s3:prop:caracter-determinante}
L'application
\[
\delta_A:(\mathbb N_0,+)\longrightarrow(\mathbb R_{>0},\cdot),
\qquad
\delta_A(k)=D_A^k,
\]
est l'unique caractère multiplicatif dont la valeur sur une prolongation
élémentaire est \(D_A\).
\end{proposition}

\begin{proof}
La multiplicativité donne
\[
\delta_A(k)
=\delta_A(\underbrace{1+\cdots+1}_{k\text{ fois}})
=\delta_A(1)^k=D_A^k.
\]
La même identité prouve l'existence et l'unicité.
\end{proof}

Le déterminant reste invariant par conjugaison de \(\mathcal T\) et par échange des feuillets \(R\mapsto-R\). La queue construite à partir de \(D_A\) appartient donc au mode commun. L'orientation apparaîtra dans le signe avec lequel le caractère régional corrige la phase.

\subsection{Règles du lecteur analytique régional}

L'extension analytique est fixée par les règles suivantes.

\begin{definition}[Lecteur régional déterminantal]
\label{p12-83-c13-s3:def:lector-regional-determinantal}
Le lecteur régional déterminantal satisfait :
\begin{enumerate}
  \item \emph{compatibilité de degré} : la longueur cylindrique est le degré
  analytique;
  \item \emph{décomposition de renouvellement} : l'impulsion finie de retour et
  la continuation stricte appartiennent à des termes additifs distincts ;
  \item \emph{normalisation de continuation} : après l'enregistrement, une seule
  fois, du résidu \(\rho_\pi\), la branche survivante commence à l'unité de
  la droite déterminant ;
  \item \emph{covariance par concaténation} : chaque prolongation élémentaire
  agit au moyen de
  \(\bigwedge^2\mathcal T=D_A\) ;
  \item \emph{orientation} : dans la convention cardinale fixée pour APP, le retour régional appartient au secteur compensateur et est soustrait de la phase d'action de cinquième degré.
\end{enumerate}
\end{definition}

La troisième règle a un contenu mathématique. La valeur \(169\) est enregistrée
une fois comme amplitude du retour ; elle n'est pas multipliée à nouveau à chaque
prolongation. La continuation stricte compte l'unique droite qui se poursuit et
la normalise par son unité. C'est la différence entre une récurrence de
renouvellement et l'itération homogène de l'impulsion initiale.

La nécessité de déclarer cette normalisation peut être prouvée. Pour tout
\(\lambda\in\mathbb R\), la famille
\[
F_\lambda(z)
=169z^6+\lambda\frac{D_Az^7}{1-D_Az}
\]
conserve le même retour fini et le même rapport déterminantal entre
coefficients consécutifs. Le catalogue et le rapport seuls ne distinguent pas
\(\lambda\). L'unité de la droite déterminant fixe
\(\lambda=1\) avant d'évaluer \(z=\alpha\). De cette manière, la normalisation
ne s'obtient pas en ajustant la valeur finale ; elle fait partie de la définition du
caractère.

\subsection{Récurrence de retour et somme close}

Nous définissons les coefficients \(\kappa_n\) par
\[
\kappa_6=169,
\qquad
\kappa_7=D_A,
\qquad
\kappa_{n+1}=D_A\kappa_n\quad(n\ge7).
\]
La première égalité est l'impulsion régionale ; les deux suivantes décrivent la
continuation normalisée.

\begin{theorem}[Réalisation analytique graduée]
\label{p12-83-c13-s3:thm:realizacion-analitica-graduada}
Le germe analytique compatible avec les règles de la
\cref{p12-83-c13-s3:def:lector-regional-determinantal} est unique et vaut
\[
\boxed{
\mathscr C_\pi(z)
=169z^6+\frac{D_Az^7}{1-D_Az}.}
\]
Ses coefficients satisfont
\[
\kappa_n=D_A^{n-6}\qquad(n\ge7).
\]
\end{theorem}

\begin{proof}
Écrivons
\[
F(z)=169z^6+\sum_{k\ge1}c_kz^{6+k}.
\]
La normalisation de continuation et la covariance déterminantale donnent
\[
c_1=D_A,
\qquad
c_{k+1}=D_Ac_k.
\]
Par récurrence, \(c_k=D_A^k\). La somme géométrique produit
\[
\sum_{k\ge1}D_A^kz^{6+k}
=\frac{D_Az^7}{1-D_Az}.
\]
Aucun coefficient libre ne subsiste.
\end{proof}

Comme \(0<\alpha<1\) et
\(0<D_A<1\),
\[
0<D_A\alpha<1.
\]
En conséquence, \(\mathscr C_\pi(\alpha)\) converge absolument et son
rayon de convergence est \(D_A^{-1}>1\). Numériquement,
\[
D_A=0.7751291192471564301888840964802\ldots,
\]
\[
D_A\alpha=0.00565639046986492673885079095469\ldots.
\]
La rapidité de cette contraction explique pourquoi la somme complète et sa
troncature de degré sept possèdent les mêmes quinze premiers chiffres dans la
carte finale.

\subsection{Le caractère d'action de 5ᵉ degré}

Le bloc central
\[
\mathcal B_c=\begin{pmatrix}7&2\\2&7\end{pmatrix}
\]
possède les valeurs propres \(9\) et \(5\). Le cycle dodécaphasique apporte le rang
\(12\) ; ses orbites de pas trois possèdent une longueur de quatre ; et le recensement complet
apporte le quotient exact
\[
\frac{104976}{1944}=54.
\]
À partir de ces invariants, on fixe le caractère d'action de cinquième degré
\[
\boxed{
H_5(\alpha,\varphi)
=\frac12\sqrt{\frac{1000\alpha}{\varphi}}-\alpha
+\frac9{16}\alpha^2-\frac59\alpha^3
+\frac7{48}\alpha^4-\frac1{54}\alpha^5.}
\]
Les quatre coefficients rationnels peuvent s'écrire
\[
\frac9{16}=\frac{\lambda_+}{4^2},
\qquad
\frac59=\frac{\lambda_-}{\lambda_+},
\qquad
\frac7{48}=\frac{\tfrac12\operatorname{tr}\mathcal B_c}{4\cdot12},
\qquad
\frac1{54}=\frac{1944}{104976},
\]
avec \((\lambda_+,\lambda_-)=(9,5)\). Ces identités certifient la
provenance des nombres utilisés. L'affectation successive de ces
invariants aux degrés deux, trois, quatre et cinq fait partie du caractère
analytique ; elle ne se déduit pas de la seule existence du catalogue fini.

Cette précision élimine une ambiguïté habituelle. Une combinaison
d'invariants de la construction est une définition interne et reproductible ; sa
naturalité exige les règles du lecteur qui spécifient l'ordre, les signes et la
normalisation. Ici, ces règles sont visibles et sont soumises aux mêmes
contrôles négatifs que le reste de la construction.

\subsection{Phase orientée et coordonnée angulaire conjuguée}

La phase de cinquième degré antérieure au retour est
\[
y_5=\frac\pi{90}\bigl(H_5+\alpha\bigr).
\]
Le caractère régional complet produit la correction
\(\mathscr C_\pi(\alpha)\). Avec l'orientation fixée dans la
\cref{p12-83-c13-s3:def:lector-regional-determinantal}, nous définissons
\[
\boxed{
y_*=\frac\pi{90}\bigl(H_5+\alpha\bigr)
-169\alpha^6
-\frac{D_A\alpha^7}{1-D_A\alpha}.}
\]
C'est la formule primaire. L'unité sexagésimale n'intervient qu'à l'application de la carte
\[
C_*^{\rm HMT}=\frac{180}{\pi}y_*.
\]

\begin{theorem}[Coordonnée angulaire conjuguée du lecteur régional déterminantal]
\label{p12-83-c13-s3:thm:contraangulo-regional}
Une fois fixés la clôture dodécaphasique de \(\alpha\),
\(\varphi=(1+\sqrt5)/2\) et les règles du lecteur régional déterminantal, la
phase \(y_*\) et sa coordonnée angulaire conjuguée sont déterminées sans introduire la valeur
recherchée. Leur évaluation est
\[
\boxed{
y_*=0.03706622646583826398493779409230638\ldots,}
\]
\[
\boxed{
C_*^{\rm HMT}
=2.12373833896864572426195138039336\ldots{}^\circ.}
\]
Donc, à quinze décimales,
\[
\boxed{C_*^{\rm HMT}=2.123738338968646^\circ.}
\]
\end{theorem}

\begin{proof}
Le \cref{p12-83-c13-s3:thm:unicidad-retorno-169} détermine \(169\). La clôture de
\(\alpha\) détermine \(A\), \(x\) et
\(D_A=e^{-100\pi\alpha/9}\). Le \cref{p12-83-c13-s3:thm:realizacion-analitica-graduada}
détermine la série complète. La formule de \(H_5\) contient uniquement
\(\alpha\), \(\varphi\) et des invariants structurels déclarés. Substituer
ces données dans l'expression de \(y_*\) donne la valeur imprimée. La valeur décimale
de \(C_*^{\rm HMT}\) n'apparaît dans aucun des membres de droite.
\end{proof}

La composante réduite d'action postérieure au retour est désormais une
sortie :
\[
\boxed{
\bar h_{\rm HMT}^{\rm ret}
=\frac{C_*^{\rm HMT}}2-\alpha
=H_5-\frac{90}{\pi}\mathscr C_\pi(\alpha).}
\]
Numériquement,
\[
\bar h_{\rm HMT}^{\rm ret}
=1.05457181691503906113369058472430\ldots.
\]
L'identité inverse n'a pas été utilisée pour introduire \(C_*^{\rm HMT}\) ; elle
s'obtient après la génération de \(y_*\).

\subsection{Troncature de degré \(7\) et borne exacte du reste}

L'expression initialement trouvée,
\[
y_7=\frac\pi{90}(H_5+\alpha)-169\alpha^6-D_A\alpha^7,
\]
est la troncature de degré sept du caractère complet. Elle produit
\[
C_7=2.12373833896864600265403827103919\ldots{}^\circ.
\]
Le reste omis n'est pas un terme indéterminé ; c'est la queue géométrique
exacte
\[
\mathscr C_\pi(\alpha)
-169\alpha^6-D_A\alpha^7
=\frac{D_A^2\alpha^8}{1-D_A\alpha}.
\]
Par conséquent,
\[
\left|C_7-C_*^{\rm HMT}\right|
=\frac{180}{\pi}\frac{D_A^2\alpha^8}{1-D_A\alpha}
=2.7839208689064583\ldots\times10^{-16}\;{}^\circ.
\]
La récurrence transforme ainsi une approximation de septième degré en une
réalisation analytique complète et fournit une borne close pour chaque
troncature ultérieure.

\subsection{Ablations du sélecteur et de la continuation}

La construction change lorsque l'une quelconque de ses composantes actives est retirée.
Les valeurs suivantes s'obtiennent sans réétalonner les autres termes :
\[
\begin{array}{@{}lc@{}}
\toprule
\text{variante}&\text{coordonnée angulaire conjuguée en degrés}\\
\midrule
\rho_\pi=168&2.1237383389772976863904487\ldots\\
\rho_\pi=169&2.1237383389686457242619514\ldots\\
\rho_\pi=170&2.1237383389599937621334541\ldots\\
\text{sans queue déterminantale}&2.1237383389686949415301675\ldots\\
D_A\text{ remplacé par }1&2.1237383389686313409965826\ldots\\
\bottomrule
\end{array}
\]
Ces ablations ne constituent pas à elles seules une estimation probabiliste.
Elles démontrent que le sélecteur équivariant et le caractère déterminantal sont
des composantes effectives de la sortie et qu'ils ne peuvent être supprimés tout en maintenant la
même réalisation.

\subsection{Reconstruction spectrale et transition hexade--octade}

La phase commune et la phase orientée déterminent les deux canaux
\[
q_-=e^{-x+y_*},
\qquad
q_+=e^{-x-y_*}.
\]
On retrouve exactement les invariants
\[
q_-q_+=D_A,
\qquad
\alpha=-\frac9{100\pi}\log D_A,
\qquad
C_*^{\rm HMT}=\frac{90}{\pi}\log\frac{q_-}{q_+}.
\]
La première égalité lit la contraction commune ; la troisième, la séparation
orientée des feuillets. L'échange \(R\mapsto-R\) permute
\(q_+\) et \(q_-\), conserve \(D_A\) et inverse le signe de \(y_*\).

Pour
\[
S_{90}(q)=\frac{q^{90}}{1-q^{270}},
\qquad
S_{120}(q)=\frac{q^{120}}{1-q^{360}},
\]
la fonctionnelle de transition hexade--octade s'écrit intégralement dans la
carte des phases :
\[
\Gamma_{6\to8}
=\frac{180}{\pi}xy_*
+12\bigl(S_{90}(q_-)-S_{90}(q_+)\bigr)
-\bigl(S_{120}(q_-)-S_{120}(q_+)\bigr).
\]
L'évaluation produite par la coordonnée \(C_*^{\rm HMT}\) du
\cref{p12-83-c13-s3:thm:contraangulo-regional} est
\[
\boxed{
\Gamma_{6\to8}
=0.274007672694459274747255260774\ldots.}
\]
Ainsi, la valeur interne reconnue en mécanique dimensionnelle comme fonctionnelle de
Barbero--Immirzi cesse de recevoir \(C_*^{\rm HMT}\) comme littéral indépendant :
toutes deux descendent du même couple de phases engendré dans ce chapitre.
L'interprétation géométrique et physique de \(\Gamma_{6\to8}\) appartient à l'œuvre
de mécanique dimensionnelle ; son évaluation mathématique appartient déjà à l'interface
précédemment construite.

\subsection{Action antérieure et postérieure au retour}

Le caractère de cinquième degré et la composante d'action postérieure au retour
sont deux valeurs distinctes :
\[
H_5
=1.05457181764615446962582182943306\ldots,
\]
\[
\bar h_{\rm HMT}^{\rm ret}
=1.05457181691503906113369058472430\ldots.
\]
La première est l'action locale avant l'incorporation du caractère régional ; la
seconde est l'action orientée qui reste après le retour. La correction
qui engendre la coordonnée angulaire conjuguée empêche de les identifier.

La comparaison métrologique doit respecter cette distinction. La première valeur donne
la mantisse
\[
2\pi H_5
=6.62607014999998792600111034989947\ldots,
\]
tandis que l'action postérieure donne
\[
2\pi\bar h_{\rm HMT}^{\rm ret}
=6.62607014540625433351074984759288\ldots.
\]
Dans le Système international, la valeur
\[
h=6.62607015\times10^{-34}\;\mathrm{J\,s}
\]
est exacte par définition \cite{BIPMSI}. C'est pourquoi,
\[
2\pi H_5-6.62607015
=-1.2073998889650101\ldots\times10^{-14},
\]
et ce n'est pas une égalité exacte. La proximité de la valeur antérieure au retour est
un contrôle numérique remarquable ; elle n'autorise pas à remplacer la constante définitoire
du SI par la composante postérieure ni à confondre les deux actions HMT\@.

Ce résultat détermine avec précision la situation mathématique. Le lecteur
régional engendre \(C_*^{\rm HMT}\) et, à partir de celui-ci, la fonctionnelle
\(\Gamma_{6\to8}\). Le même calcul sépare l'approximation de Planck
produite par \(H_5\) de l'action réduite qui intervient dans
la coordonnée angulaire conjuguée. La séparation n'affaiblit pas la construction : elle élimine une
identification incompatible avec les chiffres exacts et transforme chaque
comparaison en une affirmation réfutable de type bien défini.

\subsection{Transporteur relatif et disjonction spectrale des blocs}
\label{p12-83-c13-s3:sec:transportador-relativo-split}

Dans la base ordonnée des canaux, soit
\[
 R=\begin{pmatrix}0&1\\1&0\end{pmatrix},
 \qquad P_\pm=\frac{I\pm R}{2},
 \qquad
 B_0=7I+2R,
 \qquad B_1=\Adeg I+\Cdeg R.
\]
La décomposition scindée
\[
 xI+yR=\rho(x,y)e^{\eta(x,y)R},
 \quad \rho(x,y)=\sqrt{x^2-y^2},
 \quad \eta(x,y)=\operatorname{artanh}(y/x)
\]
construit l'unique transporteur
\[
 T_{\rm rel}
 =\frac{\sqrt{(\Adeg)^2-(\Cdeg)^2}}{\sqrt{45}}
   \exp\!\left[
   \left(\operatorname{artanh}\frac{\Cdeg}{\Adeg}
   -\operatorname{artanh}\frac27\right)R\right],
 \qquad T_{\rm rel}B_0=B_1.
\]
Dans la base spectrale, il agit par
\(9\mapsto\Adeg+\Cdeg\) et
\(5\mapsto\Adeg-\Cdeg\).

\begin{theorem}[Disjonction spectrale des réalisations locale et angulaire]
\label{p12-83-c13-s3:thm:no-go-entrelazador-bloques-split}
Les spectres
\[
 \operatorname{Spec}(B_0)=\{9,5\},\qquad
 \operatorname{Spec}(B_1)=\{\Adeg+\Cdeg,\Adeg-\Cdeg\}
\]
sont disjoints. En conséquence,
\[
 \operatorname{Hom}_{B_0,B_1}
 :=\{\Phi:\Phi B_0=B_1\Phi\}=\{0\}.
\]
\end{theorem}

\begin{proof}
Les bornes rationnelles construites pour les coordonnées angulaires donnent
\(7.29<\Adeg<7.30\) et \(2.12<\Cdeg<2.13\) ; donc,
\(9.41<\Adeg+\Cdeg<9.43\) et
\(5.16<\Adeg-\Cdeg<5.18\). Dans la base qui diagonalise \(R\), chaque entrée
d'un entrelaceur est multipliée par la différence entre une valeur propre
de \(B_0\) et une de \(B_1\) ; la disjonction impose qu'elles soient toutes nulles.
\end{proof}

Le transporteur relatif compare deux blocs au sein de l'algèbre scindée ;
les entrelaceurs dodécaphasique--Witt et d'incidence agissent entre des
représentations équivalentes dans leurs espaces de résidence propres. Cette séparation
type positivement les deux opérations et conserve leurs domaines.

\subsection{Chaîne de génération \(\mathcal F_\pi\to b_\pi\to\rho_\pi\to D_A\to C_*^{\rm HMT}\to\Gamma_{6\to8}\) : sélection régionale, prolongation déterminantale et coordonnée orientée}

La procédure complète peut être résumée par une chaîne d'applications :
\[
\begin{aligned}
\mathcal F_\pi
&\longmapsto b_\pi
\longmapsto\rho_\pi=169,\\
\alpha
&\longmapsto A
\longmapsto x
\longmapsto D_A,\\
(\rho_\pi,D_A)
&\longmapsto\mathscr C_\pi(\alpha),\\
(H_5,\mathscr C_\pi(\alpha))
&\longmapsto y_*
\longmapsto C_*^{\rm HMT},\\
(x,y_*)
&\longmapsto(q_+,q_-)
\longmapsto\Gamma_{6\to8}.
\end{aligned}
\]
La première ligne appartient au catalogue fini ; la deuxième, à la clôture commune de
\(\alpha\) ; la troisième, à la prolongation déterminantale ; la quatrième produit la
coordonnée orientée ; la cinquième la transporte vers l'interface dimensionnelle.

Dans cette organisation, \(\alpha\) mesure la contraction commune du transport et
\(C_*^{\rm HMT}\) mesure la séparation orientée de ses deux feuillets. Le retour
régional n'ajoute pas de constante externe : il sélectionne une amplitude finie et la
prolonge au moyen du déterminant déjà fixé par \(\alpha\). La carte en
degrés est ultérieure. L'objet primaire est le couple de phases \((x,y_*)\).

Une vérification exhaustive reconstruit les cinq régions directement
à partir du catalogue, parcourt les actions de
\(\mathfrak S_5\) et \(\mathfrak S_3\), vérifie la récurrence et somme la
série. L'entrée de \(\alpha\) est la sortie exacte de sa clôture
dodécaphasique, non un chiffre métrologique transcrit. La vérification maintient
séparés les théorèmes du catalogue, les règles du lecteur et la comparaison
externe ; c'est pourquoi aucune comparaison expérimentale n'intervient dans la fonction
génératrice.
